\chapter{IMPLEMENTATION}

\section{Technology Stack}
The prototype was engineered utilizing an open-source, highly portable technology stack to ensure it could be reliably deployed on standard, non-specialized hardware without licensing constraints.

\begin{table}[H]
\centering
\caption{Technology Stack for System Implementation}
\begin{tabularx}{\textwidth}{p{0.25\textwidth}p{0.28\textwidth}Y}
\toprule
\textbf{Component} & \textbf{Technology Used} & \textbf{Purpose / Function}\\
\midrule
Backend Framework & Python 3.10, Flask & Core server routing, API endpoints, thread management.\\
NLP Engine & SpaCy, Regular Expressions & Tokenization, entity extraction, structural string matching.\\
Machine Learning & Scikit-learn (TfidfVectorizer) & TF-IDF vectorization, Cosine similarity mathematical computations.\\
Document Parsers & PyPDF2, python-docx & Digital file traversing and binary text extraction.\\
OCR Fallback & Tesseract OCR, Pillow & Image preprocessing (binarization), optical character mapping.\\
Database & SQLite, SQLAlchemy & Lightweight persistence for user authentication and session logs.\\
Frontend UI & HTML5, CSS3, Bootstrap, JS & Responsive dashboard, dynamic weight sliders, real-time feedback.\\
Data Visualization & Matplotlib, Chart.js & Generating graphical Base64 score breakdowns directly in the browser.\\
\bottomrule
\end{tabularx}
\end{table}

\section{Project Folder and Module Structure}

The implementation is cleanly separated into specialized Python modules to adhere strictly to the separation of concerns software paradigm:

\begin{itemize}
    \item \textbf{app.py:} The main Flask application instance, handling HTTP REST requests, file uploads, and session management.
    \item \textbf{parser.py:} Encapsulates the PyPDF2, docx, and Tesseract routing logic, managing the file reading logic.
    \item \textbf{nlp\_engine.py:} Contains the SpaCy models, regex ontology dictionaries, and the core text normalization functions.
    \item \textbf{matcher.py:} Houses the scikit-learn TF-IDF mathematical computations and the weighted scoring formula logic.
    \item \textbf{xai.py:} Contains the Explainer class that translates the matcher.py output dictionaries into plain-English string representations.
    \item \textbf{database.py:} Manages the SQLite database schema for secure recruiter authentication.
\end{itemize}

\section{Backend Routing and Concurrency}
When an HR user navigates to the upload route, they submit a multipart form containing the target Job Description text and an array of binary file streams (resumes). The Flask route intercepts these files, saving them to a temporary \texttt{/uploads} directory. Crucially, a background processing thread is instantiated immediately. If 100 resumes are uploaded, processing them synchronously on the main thread would cause the HTTP request to timeout. The asynchronous thread iterates over each file, sequentially invoking:

\begin{center}
\texttt{parser.py} \(\rightarrow\) \texttt{nlp\_engine.py} \(\rightarrow\) \texttt{matcher.py} \(\rightarrow\) \texttt{xai.py}
\end{center}

Meanwhile, the frontend actively polls the backend status, and once complete, the final array of JSON objects is pushed to the results dashboard.

\section{Database and Authentication}
The project uses SQLite with SQLAlchemy to maintain lightweight local persistence. This supports recruiter user accounts, authentication, and session-level state without requiring an external database server. For a prototype environment, SQLite is appropriate because it is portable, self-contained, and easy to reset during testing. The authentication layer protects upload and result routes so that uploaded resumes and scoring outputs are available only to authorized users.

\section{Frontend Interface and Visualization}
The frontend is built with Bootstrap to provide a clean, modern Software-as-a-Service (SaaS) aesthetic suitable for enterprise users.

\begin{itemize}
    \item \textbf{Upload Screen:} Features drag-and-drop file inputs and dynamic HTML5 sliders, allowing recruiters to manually adjust the weighting formula prior to execution to suit specific roles, for example weighing soft skills higher for a management position.
    \item \textbf{Results Dashboard:} Displays candidates in a descending ranked table. A secondary view renders a Matplotlib/Chart.js radar chart visualizing the candidate's proficiency across the four core scoring dimensions: Technical Skills, Soft Skills, Experience, and Context.
    \item \textbf{XAI Expansion Panel:} Clicking on a specific candidate expands an accordion containing the XAI generated text, revealing the exact missing skills, matched skills, and extra bonus skills found, providing immediate context.
\end{itemize}

Figure \ref{fig:scoring-config} shows the AI scoring configuration interface. This screen is important because it demonstrates that the recruiter can control the weighted scoring logic before analysis begins. The core score is distributed across technical skills, years of experience, TF-IDF context match, and soft skills, while the extra skills bonus cap is managed separately to prevent uncontrolled score inflation.

\begin{figure}[H]
    \centering
    \IfFileExists{figures/ai_scoring_configuration.png}
    {\includegraphics[width=0.92\textwidth]{ai_scoring_configuration.png}}
    {\fbox{\parbox{0.85\textwidth}{\centering Place the scoring configuration screenshot at\\ \texttt{latex\_report/figures/ai\_scoring\_configuration.png}}}}
    \caption{AI Scoring Configuration Interface for Recruiter-Controlled Weights}
    \label{fig:scoring-config}
\end{figure}

Figure \ref{fig:xai-breakdown} shows the explainable AI output displayed on the results dashboard. The interface does not present the match score as a single opaque number; instead, it separates the score into technical skill contribution, soft skill contribution, experience contribution, context match contribution, and bonus contribution. The same panel also lists matched job-description skills, missing skills, bonus technical skills, and detected soft skills, making the ranking easier for recruiters to audit.

\begin{figure}[H]
    \centering
    \IfFileExists{figures/ai_score_breakdown.png}
    {\includegraphics[width=0.92\textwidth]{ai_score_breakdown.png}}
    {\fbox{\parbox{0.85\textwidth}{\centering Place the explainable AI screenshot at\\ \texttt{latex\_report/figures/ai\_score\_breakdown.png}}}}
    \caption{Explainable AI Score Breakdown on the Candidate Results Dashboard}
    \label{fig:xai-breakdown}
\end{figure}

\section{Known Implementation Constraints}
The current implementation utilizes SQLite, which is highly sufficient for a prototype but suffers from write-locking during high-concurrency enterprise batch operations. Furthermore, the skill ontology within \texttt{nlp\_engine.py} is currently hardcoded as a dictionary list. In a production environment, this would need to be linked to a dynamic, remotely updated graph database to keep pace with rapidly evolving industry terminology.
